% !TeX spellcheck = en
\documentclass{template} % Use the custom resume.cls style

\usepackage[left=0.5in,top=0.4in,right=0.5in,bottom=0.4in]{geometry} % Document margins
\usepackage[colorlinks=true,linkcolor=black,urlcolor=black]{hyperref}
\usepackage{marvosym}
\usepackage{graphicx}
\usepackage{contour}
\usepackage{ulem}

\normalem

\renewcommand{\ULdepth}{2.4pt}
\contourlength{0.8pt}

\newcommand{\myuline}[1]{%
  \uline{\phantom{#1}}%
  \llap{\contour{white}{#1}}%
}

\usepackage{cmap}
\usepackage[T1]{fontenc}
\usepackage{lmodern}
\input{glyphtounicode}
\pdfgentounicode=1

%----------------------------------------------------------------------------------------
%	NAME AND ADDRESS SECTION
%----------------------------------------------------------------------------------------
\name{Zain Usmani}
\address{\href{mailto:zain.usmani10@gmail.com}{zain.usmani10@gmail.com} \\
         (226) 787-6794 \\
         \href{https://www.linkedin.com/in/zain-usmani-/}{linkedin.com/in/zain-usmani-}
    }

\begin{document}

%----------------------------------------------------------------------------------------
%	EDUCATION SECTION
%----------------------------------------------------------------------------------------

\begin{rSection}{Education}
\vspace{0.1em}
  \begin{rEducationSection}{University of Windsor}
                           {2024 - Present}
                           {BASc in Electrical and Computer Engineering (Minor in CS)}
                           {GPA: 3.7/4.0}
  \end{rEducationSection}
\end{rSection}



%----------------------------------------------------------------------------------------
%	TECHNICAL EXPERIENCE SECTION
%----------------------------------------------------------------------------------------

\begin{rSection}{Experience}

\begin{rWorkSection} {Embedded Systems Developer}
                      {June 2025 -- Present}
                      {Formula Electric - Lancer Motorsport}
                      {}
  {
    \item Analyzed signal integrity and debugged physical hardware buses using oscilloscopes and logic analyzers
    \item Designed a high-performance dashboard on RPi OS Linux, bridging C++ and Linux backend services and vehicle parameters with a QT graphical interface in real-time
    \item Built automated Python scripts to validate embedded system telemetry and data transmission integrity
  } 
  \end{rWorkSection}
  
  \begin{rWorkSection}{FPGA Vision Intern}
                      {Apr 2026 -- Aug 2026}
                      {O-Hive}
                      {}
  {
    \item Engineered a real-time computer vision pipeline in C++ and deployed it to an AMD Kria KV260 embedded Linux system using a custom CMake build architecture
    \item Quantized FastDepth weights to 8 bit integer precision using Brevitas and QONNX
    \item Containerized firmware workspace via Docker to manage dependencies and ensure cross-compilation
    \item Managed Linux V4L2 drivers and Xilinx firmware to route raw video directly to hardware accelerators
    \item Designed a C++ memory fork to split and route incoming NV12 video buffers to FPGA accelerator
    \item Corrected metric distortion by designing a algorithm using physical focal lengths to fix 4.6x error margin

}
  \end{rWorkSection}

  \begin{rWorkSection}{ML and Data Intern}
                      {Jan 2025 -- Sep 2025}
                      {MicroNano Mechatronics Lab}
                      {}
  {
    \item Programmed an ESP32 Bluetooth Low Energy pipeline in C++ to stream live IMU data at 50 Hz
    \item Executed analysis on 16000 sensor samples using Pandas and NumPy to extract feature vectors
    \item Trained XGBoost classification models to distinguish human gait patterns achieving a 78.9 F1 score
    \item Co-authoring research paper detailing the machine learning performance for biomechanical classification
  } \end{rWorkSection}

\begin{rWorkSection}{Student Campus Ambassador}
                      {Dec 2024 -- June 2026}
                      {UWindsor Recruitment}
                      {}
{
  \item Guided weekly facility tours for prospective engineering students to showcase technical laboratories
  \item Represented the engineering faculty at recruitment events such as Open House and Grade 12 Day
  
} \end{rWorkSection}
\vspace{0.1em}

\end{rSection}

 
%----------------------------------------------------------------------------------------
%	PROJECTS SECTION
%----------------------------------------------------------------------------------------

\begin{rSection}{Projects}
    % Title is outside the environment to avoid the bullet

    \textbf{RouteScout Rover (McHacks 13) } $|$ \textit{Embedded C++, RPi 5, ESP32, 3D CAD} \hfill Jan 2026
            \vspace{0.2em}
    \begin{rProjectSection}
        \item Developed a real-time odometry pipeline on ESP32, which relays data back to RPi computer
        \item Captured high frequency data from motor encoders and IMU registers via hardware interrupts with I2C
        \item Implemented a Madgwick filter and closed-loop PID control to generate precise PWM motor actuation
    \end{rProjectSection}

    \vspace{0.3em} 
    \textbf{VGA Timing Controller and Frame Buffer } $|$ \textit{Verilog, Intel Quartus, MAX 10 FPGA} \hfill {Apr 2026}
        \vspace{0.2em}
    \begin{rProjectSection}
        \item Designed a hardware VGA controller in Verilog to drive a 640x480 pixel display on a MAX 10 FPGA
        \item Designed dual-port SRAM buffer to decouple video timing requirements from pixel generation logic
        \item Managed memory read timing constraints and synthesized RTL design to achieve target clock frequencies
    \end{rProjectSection}
    
\end{rSection}
\vspace{0.4em}

%----------------------------------------------------------------------------------------
%	SUMMARY / SKILLS SECTION
%----------------------------------------------------------------------------------------

\begin{rSection}{Skills}
\vspace{0.1em}
  \begin{rSummarySection}
{
\item \textbf{Languages and Tools:} C++, C, Python, Linux (Ubuntu on SoM), Shell Scripting, CMake/Ninja, Git\vspace{0.2em}
    \item \textbf{Hardware:} Oscilloscope, Waveform Generator, Logic Analyzer, Multimeters, SystemVerilog, Quartus\vspace{0.2em}
    \item \textbf{Embedded Systems:} AMD Kria KV260, STM32, ESP32, Raspberry Pi, I2C, SPI, USART, Cube IDE\vspace{0.2em}
    \item \textbf{Concepts:} Microarchitecture, Signal Integrity, I2C and SPI Protocols, DMA, PID Control, Sensor Fusion 

  }
  \end{rSummarySection}
\end{rSection}
\vspace{0.3em}


\end{document}
